\chapter{Руководство оператора}

Вкладка RS-485 предназначена для проверки функций микросервиса RS-485 через
графический интерфейс. Она позволяет отправлять данные в выбранный канал
RS-485, передавать содержимое файла и получать данные от драйвера в потоковом
режиме.

Микросервис является промежуточным слоем между графическим интерфейсом и
драйвером RS-485. Взаимодействие выполняется по следующей схеме:

\begin{center}
	\texttt{GUI --- RS-485 Microservice --- RS-485 Driver}
\end{center}

Графический интерфейс вызывает методы микросервиса по gRPC. Микросервис
обрабатывает запросы, обращается к драйверу RS-485 и возвращает результат
оператору.

\section{Подготовка к работе}

Перед использованием вкладки должны быть запущены:

\begin{itemize}
	\item драйвер RS-485 либо тестовая реализация \texttt{rs485\_fake\_driver};
	\item микросервис \texttt{rs485\_microservice};
	\item общий графический интерфейс.
\end{itemize}

Адрес микросервиса и адрес драйвера задаются в конфигурационном файле проекта
\texttt{rs/config/config.yaml}. Оператор не вводит адрес подключения во
вкладке RS-485.

При создании вкладки графический интерфейс автоматически загружает
конфигурацию и выполняет попытку подключения к микросервису. Если подключение
установлено, элементы отправки данных и управления подпиской становятся
доступными. Если подключение не установлено, эти элементы остаются
заблокированными, а причина отображается в строке состояния.

Отдельная кнопка подключения в интерфейсе отсутствует. После устранения
ошибки подключения необходимо повторно запустить графический интерфейс, чтобы
выполнить новую попытку соединения.

\section{Элементы вкладки}

\begin{longtable}{|p{0.27\textwidth}|p{0.23\textwidth}|p{0.40\textwidth}|}
	\hline
	\textbf{Элемент} & \textbf{Значение} & \textbf{Назначение} \\
	\hline
	\endfirsthead
	\hline
	\textbf{Элемент} & \textbf{Значение} & \textbf{Назначение} \\
	\hline
	\endhead
	Поле ``Канал UART'' & 0--15 & Задает идентификатор канала, используемый при вызове \texttt{SendData}. Значение по умолчанию --- 0. \\
	\hline
	Поле ``Байты'' & Последовательность байтов & Используется для ручного ввода данных в шестнадцатеричном виде. \\
	\hline
	Поле ``Файл'' & Путь к выбранному файлу & Отображает путь к файлу, содержимое которого будет отправлено вместо ручного ввода. Поле недоступно для непосредственного редактирования. \\
	\hline
	Кнопка ``Обзор...'' & --- & Открывает окно выбора файла с данными. \\
	\hline
	Кнопка ``Отправить'' & --- & Выполняет отправку вручную введенных байтов либо содержимого выбранного файла. \\
	\hline
	Кнопка ``Запустить подписку'' & --- & Открывает потоковый вызов \texttt{Subscribe}. Во время активной подписки кнопка имеет текст ``Остановить подписку''. \\
	\hline
	Поле ``Отправленные данные'' & --- & Содержит сведения о выполненных операциях отправки. \\
	\hline
	Поле ``Принятые данные'' & --- & Содержит принятые пакеты, сообщения об ошибках и события управления подпиской. \\
	\hline
	Строка состояния & --- & Показывает состояние подключения, отправки, выбора файла, подписки и получения данных. \\
	\hline
\end{longtable}

\section{Автоматическое подключение к микросервису}

Подключение выполняется автоматически во время создания вкладки. Графический
интерфейс загружает адрес микросервиса из конфигурационного файла и создает
gRPC-клиентское соединение.

До завершения попытки подключения поля канала, байтов и файла, а также кнопки
``Обзор...'', ``Отправить'' и ``Запустить подписку'' заблокированы.

При успешном подключении элементы управления разблокируются, а строка
состояния содержит сообщение:

\begin{center}
	\texttt{Статус: подключено к микросервису <адрес>}
\end{center}

При неудачном подключении элементы управления остаются заблокированными. В
строке состояния отображается адрес микросервиса либо текст возникшего
исключения.

Успешное соединение графического интерфейса с микросервисом не гарантирует
доступность драйвера RS-485. Ошибка соединения между микросервисом и драйвером
может быть получена при вызове \texttt{SendData} или во время работы
\texttt{Subscribe}.

\section{Ручной ввод данных}

Для ручной отправки данных необходимо:

\begin{enumerate}
	\item выбрать значение в поле ``Канал UART'';
	\item убедиться, что поле ``Файл'' пусто;
	\item ввести последовательность байтов в поле ``Байты'';
	\item нажать кнопку ``Отправить''.
\end{enumerate}

Байты задаются в шестнадцатеричном виде и разделяются пробелами. Каждый байт
должен находиться в диапазоне \texttt{00--FF}.

Пример корректного ввода:

\begin{center}
	\texttt{01 02 0A FF}
\end{center}

Регистр шестнадцатеричных цифр значения не имеет. Например, значения
\texttt{0a} и \texttt{0A} обозначают один и тот же байт.

Пустая строка, неполный байт, недопустимый символ или значение вне диапазона
приводят к ошибке разбора данных. Текст ошибки отображается в строке состояния
и в поле ``Отправленные данные''.

При редактировании поля ``Байты'' ранее выбранный путь к файлу автоматически
очищается. После этого источником данных становится ручной ввод.

\section{Отправка данных из файла}

Для отправки содержимого файла необходимо:

\begin{enumerate}
	\item выбрать значение в поле ``Канал UART'';
	\item нажать кнопку ``Обзор...'';
	\item выбрать требуемый файл;
	\item убедиться, что путь появился в поле ``Файл'';
	\item нажать кнопку ``Отправить''.
\end{enumerate}

После выбора файла поле ``Байты'' автоматически очищается. Если поле ``Файл''
не пусто, при отправке используется содержимое файла, а не ручной ввод.

Файл передается как последовательность байтов. В журнале отправленных данных
источник отмечается как \texttt{file: <путь>}, а данные обозначаются строкой
\texttt{binary file}. При ручном вводе источник обозначается как
\texttt{manual input}.

Если оператор закрыл окно выбора файла, не выбрав файл, текущее состояние
полей не изменяется.

\section{Выполнение SendData}

После нажатия кнопки ``Отправить'' интерфейс формирует запрос с идентификатором
канала и данными. Затем выполняется gRPC-вызов \texttt{SendData}.

В поле ``Отправленные данные'' записываются:

\begin{itemize}
	\item \texttt{channel\_id} --- идентификатор канала;
	\item \texttt{source} --- источник данных: ручной ввод или файл;
	\item \texttt{data} --- введенные байты либо обозначение двоичного файла;
	\item \texttt{success} --- признак успешного выполнения;
	\item \texttt{message} --- сообщение микросервиса;
	\item \texttt{exception} --- текст исключения, если операция не была выполнена.
\end{itemize}

При успешной отправке строка состояния содержит сообщение
``Статус: данные отправлены''. При неуспешном результате отображается сообщение
``Статус: ошибка отправки'' и текст причины.

\section{Получение данных Subscribe}

Для начала получения данных необходимо нажать кнопку
``Запустить подписку''. Интерфейс открывает потоковый gRPC-вызов
\texttt{Subscribe}, после чего:

\begin{itemize}
	\item кнопка меняет текст на ``Остановить подписку'';
	\item в поле ``Принятые данные'' появляется сообщение \texttt{Subscribe started. Waiting for data...};
	\item строка состояния сообщает о запуске подписки.
\end{itemize}

Для каждого результата, полученного из потока, в поле ``Принятые данные''
выводятся:

\begin{itemize}
	\item \texttt{channel\_id} --- идентификатор канала;
	\item \texttt{success} --- признак успешного получения;
	\item \texttt{data} --- принятые байты в шестнадцатеричном виде;
	\item \texttt{message} --- сообщение об ошибке либо пустая строка.
\end{itemize}

При успешном получении пакета строка состояния содержит сообщение
``Статус: получены данные''. При ошибке отображаются сообщение
``Статус: ошибка приема'' и текст причины.

\section{Остановка и завершение подписки}

Для остановки активной подписки оператор нажимает кнопку
``Остановить подписку''. После этого:

\begin{itemize}
	\item потоковый вызов отменяется;
	\item кнопка возвращает текст ``Запустить подписку'';
	\item в поле ``Принятые данные'' добавляется сообщение \texttt{Subscribe stopped.};
	\item строка состояния сообщает об остановке подписки.
\end{itemize}

Поток также может завершиться без нажатия кнопки, например при разрыве
соединения или остановке драйвера. В этом случае кнопка автоматически
возвращается в исходное состояние, а в поле ``Принятые данные'' добавляется
сообщение \texttt{Subscribe finished.}.

Если перед завершением потока была показана ошибка приема, она сохраняется в
строке состояния. В остальных случаях строка состояния сообщает:
``Статус: подписка завершена''.

При закрытии графического интерфейса клиент отключается от микросервиса, а
активная подписка завершается.

\section{Коды ошибок микросервиса}

\begin{longtable}{|p{0.35\textwidth}|p{0.55\textwidth}|}
	\hline
	\textbf{Код} & \textbf{Описание} \\
	\hline
	\endfirsthead
	\hline
	\textbf{Код} & \textbf{Описание} \\
	\hline
	\endhead
	\texttt{NO\_ERROR} & Операция выполнена успешно. \\
	\hline
	\texttt{INVALID\_CHANNEL\_ID} & Передан недопустимый идентификатор канала. \\
	\hline
	\texttt{INVALID\_DATA} & Данные отсутствуют либо имеют недопустимый формат. \\
	\hline
	\texttt{DRIVER\_NOT\_CONNECTED} & Микросервис не подключен к драйверу RS-485. \\
	\hline
	\texttt{DRIVER\_NO\_REPLY} & Драйвер или подключенное устройство не вернуло ответ. \\
	\hline
	\texttt{DRIVER\_TIMEOUT} & Истекло время ожидания ответа. \\
	\hline
	\texttt{DRIVER\_ERROR} & Драйвер сообщил об ошибке выполнения операции. \\
	\hline
	\texttt{INTERNAL\_ERROR} & Произошла внутренняя ошибка микросервиса или gRPC-взаимодействия. \\
	\hline
\end{longtable}

\section{Действия при ошибках}

Если элементы вкладки остаются заблокированными, необходимо:

\begin{enumerate}
	\item убедиться, что микросервис RS-485 запущен;
	\item проверить адрес сервиса в файле \texttt{rs/config/config.yaml};
	\item убедиться, что указанный TCP-порт свободен;
	\item проверить доступность конфигурационного файла для GUI;
	\item повторно запустить графический интерфейс.
\end{enumerate}

Если подключение к микросервису установлено, но отправка или подписка
завершаются ошибкой, необходимо проверить:

\begin{itemize}
	\item корректность идентификатора канала;
	\item формат вручную введенных байтов;
	\item существование и доступность выбранного файла;
	\item работу драйвера RS-485 либо \texttt{rs485\_fake\_driver};
	\item адрес драйвера в конфигурационном файле;
	\item физическое соединение RS-485 при использовании реального устройства.
\end{itemize}

При ошибках \texttt{DRIVER\_NO\_REPLY} и \texttt{DRIVER\_TIMEOUT} необходимо
дополнительно проверить доступность подключенного устройства.

При ошибке \texttt{INTERNAL\_ERROR} рекомендуется сохранить текст сообщения
из строки состояния и журналов и передать его разработчику.

\section{Завершение работы}

Перед закрытием графического интерфейса рекомендуется остановить активную
подписку. После этого можно закрыть GUI, остановить микросервис и затем
остановить драйвер RS-485.

Для остановки программ, запущенных из терминала, используется сочетание клавиш
\texttt{Ctrl+C}.